% Options for packages loaded elsewhere
% Options for packages loaded elsewhere
\PassOptionsToPackage{unicode}{hyperref}
\PassOptionsToPackage{hyphens}{url}
\PassOptionsToPackage{dvipsnames,svgnames,x11names}{xcolor}
%
\documentclass[
  letterpaper,
  DIV=11,
  numbers=noendperiod]{scrartcl}
\usepackage{xcolor}
\usepackage{amsmath,amssymb}
\setcounter{secnumdepth}{-\maxdimen} % remove section numbering
\usepackage{iftex}
\ifPDFTeX
  \usepackage[T1]{fontenc}
  \usepackage[utf8]{inputenc}
  \usepackage{textcomp} % provide euro and other symbols
\else % if luatex or xetex
  \usepackage{unicode-math} % this also loads fontspec
  \defaultfontfeatures{Scale=MatchLowercase}
  \defaultfontfeatures[\rmfamily]{Ligatures=TeX,Scale=1}
\fi
\usepackage{lmodern}
\ifPDFTeX\else
  % xetex/luatex font selection
\fi
% Use upquote if available, for straight quotes in verbatim environments
\IfFileExists{upquote.sty}{\usepackage{upquote}}{}
\IfFileExists{microtype.sty}{% use microtype if available
  \usepackage[]{microtype}
  \UseMicrotypeSet[protrusion]{basicmath} % disable protrusion for tt fonts
}{}
\makeatletter
\@ifundefined{KOMAClassName}{% if non-KOMA class
  \IfFileExists{parskip.sty}{%
    \usepackage{parskip}
  }{% else
    \setlength{\parindent}{0pt}
    \setlength{\parskip}{6pt plus 2pt minus 1pt}}
}{% if KOMA class
  \KOMAoptions{parskip=half}}
\makeatother
% Make \paragraph and \subparagraph free-standing
\makeatletter
\ifx\paragraph\undefined\else
  \let\oldparagraph\paragraph
  \renewcommand{\paragraph}{
    \@ifstar
      \xxxParagraphStar
      \xxxParagraphNoStar
  }
  \newcommand{\xxxParagraphStar}[1]{\oldparagraph*{#1}\mbox{}}
  \newcommand{\xxxParagraphNoStar}[1]{\oldparagraph{#1}\mbox{}}
\fi
\ifx\subparagraph\undefined\else
  \let\oldsubparagraph\subparagraph
  \renewcommand{\subparagraph}{
    \@ifstar
      \xxxSubParagraphStar
      \xxxSubParagraphNoStar
  }
  \newcommand{\xxxSubParagraphStar}[1]{\oldsubparagraph*{#1}\mbox{}}
  \newcommand{\xxxSubParagraphNoStar}[1]{\oldsubparagraph{#1}\mbox{}}
\fi
\makeatother

\usepackage{color}
\usepackage{fancyvrb}
\newcommand{\VerbBar}{|}
\newcommand{\VERB}{\Verb[commandchars=\\\{\}]}
\DefineVerbatimEnvironment{Highlighting}{Verbatim}{commandchars=\\\{\}}
% Add ',fontsize=\small' for more characters per line
\usepackage{framed}
\definecolor{shadecolor}{RGB}{241,243,245}
\newenvironment{Shaded}{\begin{snugshade}}{\end{snugshade}}
\newcommand{\AlertTok}[1]{\textcolor[rgb]{0.68,0.00,0.00}{#1}}
\newcommand{\AnnotationTok}[1]{\textcolor[rgb]{0.37,0.37,0.37}{#1}}
\newcommand{\AttributeTok}[1]{\textcolor[rgb]{0.40,0.45,0.13}{#1}}
\newcommand{\BaseNTok}[1]{\textcolor[rgb]{0.68,0.00,0.00}{#1}}
\newcommand{\BuiltInTok}[1]{\textcolor[rgb]{0.00,0.23,0.31}{#1}}
\newcommand{\CharTok}[1]{\textcolor[rgb]{0.13,0.47,0.30}{#1}}
\newcommand{\CommentTok}[1]{\textcolor[rgb]{0.37,0.37,0.37}{#1}}
\newcommand{\CommentVarTok}[1]{\textcolor[rgb]{0.37,0.37,0.37}{\textit{#1}}}
\newcommand{\ConstantTok}[1]{\textcolor[rgb]{0.56,0.35,0.01}{#1}}
\newcommand{\ControlFlowTok}[1]{\textcolor[rgb]{0.00,0.23,0.31}{\textbf{#1}}}
\newcommand{\DataTypeTok}[1]{\textcolor[rgb]{0.68,0.00,0.00}{#1}}
\newcommand{\DecValTok}[1]{\textcolor[rgb]{0.68,0.00,0.00}{#1}}
\newcommand{\DocumentationTok}[1]{\textcolor[rgb]{0.37,0.37,0.37}{\textit{#1}}}
\newcommand{\ErrorTok}[1]{\textcolor[rgb]{0.68,0.00,0.00}{#1}}
\newcommand{\ExtensionTok}[1]{\textcolor[rgb]{0.00,0.23,0.31}{#1}}
\newcommand{\FloatTok}[1]{\textcolor[rgb]{0.68,0.00,0.00}{#1}}
\newcommand{\FunctionTok}[1]{\textcolor[rgb]{0.28,0.35,0.67}{#1}}
\newcommand{\ImportTok}[1]{\textcolor[rgb]{0.00,0.46,0.62}{#1}}
\newcommand{\InformationTok}[1]{\textcolor[rgb]{0.37,0.37,0.37}{#1}}
\newcommand{\KeywordTok}[1]{\textcolor[rgb]{0.00,0.23,0.31}{\textbf{#1}}}
\newcommand{\NormalTok}[1]{\textcolor[rgb]{0.00,0.23,0.31}{#1}}
\newcommand{\OperatorTok}[1]{\textcolor[rgb]{0.37,0.37,0.37}{#1}}
\newcommand{\OtherTok}[1]{\textcolor[rgb]{0.00,0.23,0.31}{#1}}
\newcommand{\PreprocessorTok}[1]{\textcolor[rgb]{0.68,0.00,0.00}{#1}}
\newcommand{\RegionMarkerTok}[1]{\textcolor[rgb]{0.00,0.23,0.31}{#1}}
\newcommand{\SpecialCharTok}[1]{\textcolor[rgb]{0.37,0.37,0.37}{#1}}
\newcommand{\SpecialStringTok}[1]{\textcolor[rgb]{0.13,0.47,0.30}{#1}}
\newcommand{\StringTok}[1]{\textcolor[rgb]{0.13,0.47,0.30}{#1}}
\newcommand{\VariableTok}[1]{\textcolor[rgb]{0.07,0.07,0.07}{#1}}
\newcommand{\VerbatimStringTok}[1]{\textcolor[rgb]{0.13,0.47,0.30}{#1}}
\newcommand{\WarningTok}[1]{\textcolor[rgb]{0.37,0.37,0.37}{\textit{#1}}}

\usepackage{longtable,booktabs,array}
\newcounter{none} % for unnumbered tables
\usepackage{calc} % for calculating minipage widths
% Correct order of tables after \paragraph or \subparagraph
\usepackage{etoolbox}
\makeatletter
\patchcmd\longtable{\par}{\if@noskipsec\mbox{}\fi\par}{}{}
\makeatother
% Allow footnotes in longtable head/foot
\IfFileExists{footnotehyper.sty}{\usepackage{footnotehyper}}{\usepackage{footnote}}
\makesavenoteenv{longtable}
\usepackage{graphicx}
\makeatletter
\newsavebox\pandoc@box
\newcommand*\pandocbounded[1]{% scales image to fit in text height/width
  \sbox\pandoc@box{#1}%
  \Gscale@div\@tempa{\textheight}{\dimexpr\ht\pandoc@box+\dp\pandoc@box\relax}%
  \Gscale@div\@tempb{\linewidth}{\wd\pandoc@box}%
  \ifdim\@tempb\p@<\@tempa\p@\let\@tempa\@tempb\fi% select the smaller of both
  \ifdim\@tempa\p@<\p@\scalebox{\@tempa}{\usebox\pandoc@box}%
  \else\usebox{\pandoc@box}%
  \fi%
}
% Set default figure placement to htbp
\def\fps@figure{htbp}
\makeatother





\setlength{\emergencystretch}{3em} % prevent overfull lines

\providecommand{\tightlist}{%
  \setlength{\itemsep}{0pt}\setlength{\parskip}{0pt}}



 


% Page furniture for the activity PDFs, ported from the Physics 230A repo, which in turn ported it from the exam-documentclass originals that course used before its Quarto migration.
%
% Pulled in by every problem set's YAML:
%
%   format:
%     pdf:
%       include-in-header: header.tex
%     latex:
%       include-in-header: header.tex
%
% Both keys are needed: pdf and latex are separate formats to Quarto, and with the key on pdf alone the shipped .tex comes out without the header its own .pdf has.
%
% The exam class provided \header, \headrule and \footer; Quarto's pdf output uses the default article class, so fancyhdr supplies the same furniture.
% The exam class's \numpages becomes \pageref*{LastPage}, which needs two passes --- Quarto runs latexmk, so that is already handled.
% The star matters: plain \pageref is a hyperref link, so the page count comes out link-blue.
%
% ---------------------------------------------------------------------------
% ADAPTING THIS FOR ANOTHER COURSE: change the two strings below, nothing else.
% ---------------------------------------------------------------------------
%
% Everything below is commented out, as it is in the Physics 230A repo this came from: the wiring is in place and the furniture is off.
% Uncommenting the block turns on a running header and footer in every activity PDF, and in the syllabus PDF if this repo ever grows one.

% \usepackage{fancyhdr}
% \usepackage{lastpage}

% \pagestyle{fancy}
% \fancyhf{}

% \fancyhead[L]{Modern DevOps for Systems Biologists}  % <- course name
% \fancyfoot[L]{\textit{jun.allard@uci.edu}}           % <- contact

% \fancyfoot[R]{Page \thepage\ of \pageref*{LastPage}}

% \renewcommand{\headrulewidth}{0.4pt}
% \renewcommand{\footrulewidth}{0pt}

% % Quarto's title block starts the document with \maketitle, which forces the plain pagestyle on the first page and would otherwise drop the header exactly where it is most wanted.
% \fancypagestyle{plain}{%
%   \fancyhf{}%
%   \fancyhead[L]{Modern DevOps for Systems Biologists}%
%   \fancyfoot[L]{\textit{jun.allard@uci.edu}}%
%   \fancyfoot[R]{Page \thepage\ of \pageref*{LastPage}}%
%   \renewcommand{\headrulewidth}{0.4pt}%
%   \renewcommand{\footrulewidth}{0pt}%
% }
\KOMAoption{captions}{tableheading}
\makeatletter
\@ifpackageloaded{caption}{}{\usepackage{caption}}
\AtBeginDocument{%
\ifdefined\contentsname
  \renewcommand*\contentsname{Table of contents}
\else
  \newcommand\contentsname{Table of contents}
\fi
\ifdefined\listfigurename
  \renewcommand*\listfigurename{List of Figures}
\else
  \newcommand\listfigurename{List of Figures}
\fi
\ifdefined\listtablename
  \renewcommand*\listtablename{List of Tables}
\else
  \newcommand\listtablename{List of Tables}
\fi
\ifdefined\figurename
  \renewcommand*\figurename{Figure}
\else
  \newcommand\figurename{Figure}
\fi
\ifdefined\tablename
  \renewcommand*\tablename{Table}
\else
  \newcommand\tablename{Table}
\fi
}
\@ifpackageloaded{float}{}{\usepackage{float}}
\floatstyle{ruled}
\@ifundefined{c@chapter}{\newfloat{codelisting}{h}{lop}}{\newfloat{codelisting}{h}{lop}[chapter]}
\floatname{codelisting}{Listing}
\newcommand*\listoflistings{\listof{codelisting}{List of Listings}}
\makeatother
\makeatletter
\makeatother
\makeatletter
\@ifpackageloaded{caption}{}{\usepackage{caption}}
\@ifpackageloaded{subcaption}{}{\usepackage{subcaption}}
\makeatother
\usepackage{bookmark}
\IfFileExists{xurl.sty}{\usepackage{xurl}}{} % add URL line breaks if available
\urlstyle{same}
\hypersetup{
  pdftitle={Activity 4: Two virtual environments},
  pdfauthor={Jun Allard},
  colorlinks=true,
  linkcolor={blue},
  filecolor={Maroon},
  citecolor={Blue},
  urlcolor={Blue},
  pdfcreator={LaTeX via pandoc}}

\ifXeTeX
\special{pdf:trailerid [ <f61c361cf8716d6f1e8bec6559f0dc05> <f61c361cf8716d6f1e8bec6559f0dc05> ]}
\fi
\ifPDFTeX
\pdftrailerid{}
\pdftrailer{/ID [ <f61c361cf8716d6f1e8bec6559f0dc05> <f61c361cf8716d6f1e8bec6559f0dc05> ]}
\fi
\ifLuaTeX
\pdfvariable trailerid {[ <f61c361cf8716d6f1e8bec6559f0dc05> <f61c361cf8716d6f1e8bec6559f0dc05> ]}
\fi


\title{Activity 4: Two virtual environments}
\usepackage{etoolbox}
\makeatletter
\providecommand{\subtitle}[1]{% add subtitle to \maketitle
  \apptocmd{\@title}{\par {\large #1 \par}}{}{}
}
\makeatother
\subtitle{Modern DevOps for Systems Biologists}
\author{Jun Allard}
\date{}
\begin{document}
\maketitle


Last year you measured the brightest pixel in a microscope image, added
the background offset your protocol calls for, and wrote the number in
your paper. Today you run the same script on the same image, and get a
different number.

\subsection{A. Two answers from one
script}\label{a.-two-answers-from-one-script}

\begin{enumerate}
\def\labelenumi{\roman{enumi}.}
\item
  Check that you have \texttt{uv}, the tool that builds and runs Python
  environments:

\begin{Shaded}
\begin{Highlighting}[]
\ExtensionTok{uv} \AttributeTok{{-}{-}version}
\end{Highlighting}
\end{Shaded}

  If that fails, install it and open a new terminal:

\begin{Shaded}
\begin{Highlighting}[]
\ExtensionTok{curl} \AttributeTok{{-}LsSf}\NormalTok{ https://astral.sh/uv/install.sh }\KeywordTok{|} \FunctionTok{sh}
\end{Highlighting}
\end{Shaded}
\item
  Write this into a file called \texttt{brightest.py}:

\begin{Shaded}
\begin{Highlighting}[]
\ImportTok{import}\NormalTok{ numpy }\ImportTok{as}\NormalTok{ np}

\CommentTok{\# An 8{-}bit microscope image. Every pixel is a whole number from 0 to 255.}
\NormalTok{image }\OperatorTok{=}\NormalTok{ np.array([[}\DecValTok{12}\NormalTok{, }\DecValTok{250}\NormalTok{], [}\DecValTok{180}\NormalTok{, }\DecValTok{7}\NormalTok{]], dtype}\OperatorTok{=}\NormalTok{np.uint8)}

\NormalTok{peak }\OperatorTok{=}\NormalTok{ image.}\BuiltInTok{max}\NormalTok{()      }\CommentTok{\# the brightest pixel}
\NormalTok{corrected }\OperatorTok{=}\NormalTok{ peak }\OperatorTok{+} \DecValTok{100}  \CommentTok{\# the protocol says to add a background offset of 100}

\BuiltInTok{print}\NormalTok{(}\StringTok{"numpy"}\NormalTok{, np.\_\_version\_\_)}
\BuiltInTok{print}\NormalTok{(}\StringTok{"brightest pixel:"}\NormalTok{, peak)}
\BuiltInTok{print}\NormalTok{(}\StringTok{"after offset:   "}\NormalTok{, corrected)}
\end{Highlighting}
\end{Shaded}
\item
  Run it twice, against two different versions of numpy.
  \texttt{uv\ run\ -\/-with} builds a throwaway environment containing
  exactly the version you name, runs your script inside it, and throws
  it away again:

\begin{Shaded}
\begin{Highlighting}[]
\ExtensionTok{uv}\NormalTok{ run }\AttributeTok{{-}{-}with} \StringTok{\textquotesingle{}numpy==1.26.4\textquotesingle{}}\NormalTok{ brightest.py}
\ExtensionTok{uv}\NormalTok{ run }\AttributeTok{{-}{-}with} \StringTok{\textquotesingle{}numpy==2.3.2\textquotesingle{}}\NormalTok{  brightest.py}
\end{Highlighting}
\end{Shaded}
\item
  Write down both answers.

  (If \texttt{uv} pauses to download a Python as well as a numpy, that
  is expected: the interpreter is part of the environment too, and
  \texttt{uv} picks one each version of numpy can actually run on.)
\end{enumerate}

\subsection{B. Which number is wrong, and
why}\label{b.-which-number-is-wrong-and-why}

\begin{enumerate}
\def\labelenumi{\roman{enumi}.}
\item
  The brightest pixel is 250 in both runs, so the image is not the
  problem. 250 + 100 = 350, and 350 does not fit in the 0-255 range that
  \texttt{uint8} can hold. One version widens the number to make it fit;
  the other lets it wrap around past 255 and start again from 0. Say
  which version did which, and check that the wrapped answer is the one
  you would predict.
\item
  Only one of the two runs said anything about this at all. Look again
  at the output of each --- including anything printed before the
  numbers --- and say what warning you got, from which version, and
  whether you would have noticed it in a script that printed a hundred
  lines.
\item
  This change is documented in
  \href{https://numpy.org/neps/nep-0050-scalar-promotion.html}{NEP 50},
  which shipped in numpy 2.0.
\end{enumerate}

\subsection{C. Say which version you
meant}\label{c.-say-which-version-you-meant}

Make the version a property of the project.

\begin{enumerate}
\def\labelenumi{\roman{enumi}.}
\item
  Create a project and pin the version you decided to believe:

\begin{Shaded}
\begin{Highlighting}[]
\ExtensionTok{uv}\NormalTok{ init brightness}
\BuiltInTok{cd}\NormalTok{ brightness}
\ExtensionTok{uv}\NormalTok{ add }\StringTok{\textquotesingle{}numpy==1.26.4\textquotesingle{}}      \CommentTok{\# or 2.3.2 {-}{-}{-} your choice, but say which}
\end{Highlighting}
\end{Shaded}
\item
  Move \texttt{brightest.py} into that directory and run it with no
  \texttt{-\/-with} at all:

\begin{Shaded}
\begin{Highlighting}[]
\ExtensionTok{uv}\NormalTok{ run brightest.py}
\end{Highlighting}
\end{Shaded}

  Confirm you get the answer you chose.
\item
  Two files now record the environment.

\begin{Shaded}
\begin{Highlighting}[]
\FunctionTok{cat}\NormalTok{ pyproject.toml}
\FunctionTok{cat}\NormalTok{ uv.lock}
\end{Highlighting}
\end{Shaded}

  \texttt{pyproject.toml} is what you asked for; \texttt{uv.lock} is
  what you got, down to the exact build of every package numpy itself
  depends on. The second is the one that makes someone else's run
  identical to yours. \texttt{uv\ init} wrote a third,
  \texttt{.python-version}, for the same reason.
\item
  Commit both, together with the script. A collaborator who clones this
  and types \texttt{uv\ run\ brightest.py} now gets your number, on
  their machine, without asking you anything.
\item
  Finally, write a \texttt{QUICKSTART.md} in that directory: four lines
  that take a stranger from a fresh clone to the number your script
  prints. Assume they have nothing installed but \texttt{git}.
\end{enumerate}




\end{document}
